% !TEX root = ../../../main.tex
% 所属：第一部分：单变量分析学 / 第3章：实数的八大基本定理
% 节：3.11 聚点定理
% 正文按原文件顺序保存；由 ../chapter3.tex 调用。
\section{聚点定理}
在\ref{defoflimpoint}中已经给出了聚点的定义
\begin{theorem}
	若$E$是一个有界无限集，则$E$有聚点
\end{theorem}
\begin{proof}
	$E$ 为有界无限点集，则必有可列子集 $\{x_1, \cdots, x_n, \cdots\}$，那么 $\{x_n\}$ 有界，由Bolzano-Weierstrass致密性定理，有收敛子列 $\{x_{n_k}\}$，且 $x_{n_k}$ 趋近于某点 $a$.
	
	下证 $a$ 为聚点，由 $x_{n_k} \to a$，$\forall \delta > 0, \exists K > 0$，使得 $\forall k > K$，有
	\[
	|x_{n_k} - a| < \delta
	\]
	
	可以找到一点，属于 $a$ 的 $\delta$ 去心邻域内
\end{proof}
\begin{proposition}
	$$\text{聚点定理}\Rightarrow\text{Bolzano-Weierstrass致密性定理}$$
\end{proposition}
\begin{proof}
	$\{x_1, \cdots, x_n, \cdots\}$是有界无穷点集，有聚点$a$,下证$a$为$\{x_n\}$的某一个子列的极限
	\[\begin{aligned}
		&\delta_1=\dfrac{1}{2},\forall x_{n_1}\in\mathring{O}(a,\dfrac{1}{2})\cap E,\textbf{有}|x_{n_1}-a|<\delta_1\\
		&\delta_2=\dfrac{1}{3},\forall x_{n_2}\in\mathring{O}(a,\dfrac{1}{3})\cap E,\textbf{有}|x_{n_2}-a|<\delta_2\\
		&\vdots\\
		&\delta_1k=\dfrac{1}{k+1},\forall x_{n_k}\in\mathring{O}(a,\dfrac{1}{k+1})\cap E,\textbf{有}|x_{n_k}-a|<\delta_k
	\end{aligned}\]
	这样$x_{ n_k}\to a$
\end{proof}
